\documentclass[12pt]{report}
\usepackage{amsmath}
\usepackage{amssymb}
\usepackage{amsthm}
\usepackage{enumitem}
\usepackage{multicol}

\begin{document}
\renewcommand\qedsymbol{OE$\Delta$}
\newcommand{\abs}[1]{\left| #1 \right|}
\newcommand{\andtxt}{\text{ and }}
\newcommand{\mathif}{\text{if }}
\newcommand{\seq}[2][n]{\ensuremath{(#2 _{#1})}}

\tableofcontents
\pagebreak

\setcounter{chapter}{1}
\chapter{Theories of Cantor (Somehow not Jewish)}
\section{Exercise 2.1}
\begin{itemize}{}
    \item (a) List all the elements of $\mathcal{P}(\{a,b,c\})$. 
    \item (b) Define a formula for the number of elements in the power set of an n-element set.
\end{itemize}

\subsection{Solution (a)}
The power set is the set of all subsets.
Below is a list of all such subsets.
\begin{center}
    \begin{tabular}{|c|c|c|}
        \hline
        1 & 2 & 3 \\
        \hline
        \{a\} & \{b\} & \{c\} \\
        \hline \hline
        4 & 5 & 6 \\
        \hline
        \{a,b\} & \{a,c\} & \{b,c\} \\
        \hline \hline
        7 & 8 & \\
        \hline
        \{a,b,c\} & $\emptyset$ & \\
        \hline
    \end{tabular}
\end{center}

There are \boxed{8} members of $\mathcal{P}(\{a,b,c\})$. 

\subsection{Solution (b)}
We can look at the first couple values.
\begin{center}
    \begin{tabular}{|c|c|c|c|c|c}
        \hline
        n & 0 & 1 & 2 & 3 \\
        \hline
        $\mathcal{P}$ & 1 & 2 & 4 & 8 \\
        \hline
    \end{tabular}
\end{center}

As it stands, it appears that the proper formula for it is \boxed{2^n}. 

\pagebreak
\section{Exercise 2.2}
Prove that $\left|\left\{ e^n : n \in \mathbb{N} \right\}\right| = \left| \mathbb{N} \right|$.

\subsection*{Solution}
\begin{proof}
    Let there be a function $f$.
    \begin{gather}
        f: \mathbb{N} \rightarrow e^\mathbb{N}\\
        n \rightarrowtail e^n
    \end{gather}
    
    This pairs every natural number with a number in the left set.
    Conclusively, the two sets have the same cardinality and there is a bijection.
    THis can be written in math terms.
    \begin{equation}
        \left|\left\{ e^n : n \in \mathbb{N} \right\}\right| = \left| \mathbb{N} \right|
    \end{equation}
    Q.E.D.
\end{proof}

\pagebreak
\section{Exercise 2.3}
The following pairs of sets have the same size, and  so there exists a bijection between them. Write down an explicit  bijection in each case. You do not need to prove your answers.  
\begin{multicols}{2}
    (a) $(0, \infty)$ and $(1, \infty)$ 

    (b) $(0, \infty)$ and $(-\infty, 3)$

    (c) $(0, \infty)$ and (0, 1) 

    (d) $\mathbb{R}$ and $(0, \infty)$  
    \columnbreak

    (e) $\mathbb{R}$ and (0, 1) 
    
    (f) $\mathbb{Z}$ and $\{ ..., \frac{1}{8}, \frac{1}{4}, \frac{1}{2}, 1, 2, 4, 8, ... \}$  
    
    (g) $\{0, 1\} \times \mathbb{N}$ and $\mathbb{N}$  
    
    (h) [0, 1] and (0, 1) 
\end{multicols}

\subsection{Solution (a)}
\begin{multicols}{2}
    \begin{gather}
        f : (0,\infty) \to (1,\infty)\\
        n \rightarrowtail n + 1
    \end{gather}
    
    \columnbreak
    \begin{gather}
        f : (1,\infty) \to (0,\infty)\\
        n \rightarrowtail n - 1
    \end{gather}
\end{multicols}

\subsection{Solution (b)}
\begin{multicols}{2}
    \begin{gather}
        f : (0,\infty) \to (-\infty, 3)\\
        n \rightarrowtail -n + 3
    \end{gather}

    \columnbreak
    \begin{gather}
        f : (-\infty, 3) \to (0,\infty)\\
        n \rightarrowtail -n + 3
    \end{gather}
    
\end{multicols}

\subsection{Solution (c)}
\begin{multicols}{2}
    \begin{gather}
        f : (0,\infty) \to (0,1)\\
        n \rightarrowtail \frac{1}{n + 1}
    \end{gather}

    \columnbreak
    \begin{gather}
        f : (0,1) \to (0,\infty)\\
        n \rightarrowtail \frac{1}{n} - 1
    \end{gather}
    
\end{multicols}

\subsection{Solution (d)}
\begin{multicols}{2}
    \begin{gather}
        f : \mathbb{R} \to (0, \infty)\\
        n \rightarrowtail e^x
    \end{gather}

    \columnbreak
    \begin{gather}
        f : (0, \infty) \to \mathbb{R}\\
        n \rightarrowtail \ln(x)
    \end{gather}
\end{multicols}

\subsection{Solution (e)}
\begin{multicols}{2}
    \begin{gather}
        f : \mathbb{R} \to (0, 1)\\
        n \rightarrowtail \frac{1}{2n - 1}
    \end{gather}

    \columnbreak
    \begin{gather}
        f : (0, 1) \to \mathbb{R}\\
        n \rightarrowtail \frac{1}{2n} + \frac{1}{2}
    \end{gather}
\end{multicols}

\subsection{Solution (f)}
\begin{multicols}{2}
    \begin{gather}
        f : \mathbb{Z} \to \{ ..., \frac{1}{8}, \frac{1}{4}, \frac{1}{2}, 1, 2, 4, 8, ... \}\\
        n \rightarrowtail 2^n
    \end{gather}

    \columnbreak
    \begin{gather}
        f : \{ ..., \frac{1}{8}, \frac{1}{4}, \frac{1}{2}, 1, 2, 4, 8, ... \} \to \mathbb{Z}\\
        n \rightarrowtail \log_2 (n)
    \end{gather}
\end{multicols}

\subsection{Solution (g)}
\begin{multicols}{2}
    \begin{gather}
        f : \{0, 1\} \times \mathbb{N} \to \mathbb{N}\\
        (a,b) \rightarrowtail a \cdot b
    \end{gather}
    \columnbreak
    \begin{gather}
        f : \mathbb{N} \to \{0, 1\} \times \mathbb{N}\\
        n \rightarrowtail \left(n \mod 2, \frac{n - (n \mod 2)}{2}\right)
    \end{gather}
\end{multicols}

\subsection{Solution (h)}
Iffy solution. 
\begin{multicols}{2}
    \begin{gather}
        f : [0, 1] \to (0, 1)\\
        n \rightarrowtail \begin{cases}
            \frac{1}{2} & \text{if } n=0\\
            \frac{1}{2^{x+2}} & \text{if } n=\frac{1}{2^x}\\
            n    &   \text{otherwise}
        \end{cases}
    \end{gather}

    \columnbreak
    \begin{gather}
        f : (0, 1) \to [0,1]\\
        n \rightarrowtail \begin{cases}
            0 & \text{if } n=\frac{1}{2}\\
            \frac{1}{2^{x-2}} & \text{if } n=\frac{1}{2^x}\\
            n    &   \text{otherwise}
        \end{cases}
    \end{gather}
\end{multicols}

\pagebreak
\section{Exercise 2.4}
This problem shows that ``equinumerosity is an equivalence relation.'' (This justifies the notation $|A| = |B|$.) Let $A$, $B$, and $C$ be sets. For this problem only, we'll write $A \sim B$ to mean that $A$ and $B$ are equinumerous, meaning that there is a bijection $A \to B$. 

(a) Show that $A \sim A$.  

(b) Show that if $A \sim B$ then $B \sim A$.  

(c) Show that if $A \sim B$ and $B \sim C$, then $A \sim C$. 

\subsection{Solution (a)}
We can always establish a bijection between $A$ and $A$.
\begin{gather}
    f: A \to A \\
    n \rightarrowtail n
\end{gather}
This establishes the bijection. 

\subsection{Solution (b)}
It is given that $A \sim B$. 
We can extract from that that the function from $A$ to $B$ is bijective, and consequently that it is reversable.
Said reversability can give us the function from $B$ to $A$, which would itself be bijective. 
This bijectivity says that $B \sim A$. 

\subsection{Solution (c)}
Since $A \sim B$, then $A \to B$ is bijective.
The same can be done to conclude that $B \to C$ is bijective.
Since both are bijective, this means that $f(a) = f(b) \implies a = b$ and $f(b) = f(c) \implies b = c$, as well as that $b \in B \implies \exists a \in A : f(a) = b$ and $c \in C \implies \exists b \in B : f(b) = c$.
By transistivity, since both $A \to B$ and $B \to C$ are bijective, then $A \to C$ is bijective.
This in turn means that $A \sim C$.

\pagebreak
\section{Exercise 2.5}
\begin{itemize}
    \item (a) Prove that if A and B are countable sets, then $A \cup B$ is also a countable set.  
    
    \item (b) Prove that if $A_n$ is a countable set for each $n \in N$, then the set  $\overunderset{\infty}{n = 1}{\bigcup} A_n$ is also countable. 
\end{itemize}

\subsection{Solution (a)}
We can establish a bijective function. 
\begin{multicols}{2}
    \begin{gather}
        f: \mathbb{N} \to A \cup B\\
        n \to \begin{cases}
            A_{\lfloor\frac{x}{2}\rfloor} & \text{if } x \equiv 0 (\mod 2)\\
            B_{\lfloor\frac{x}{2}\rfloor} & \text{if } x \equiv 1 (\mod 2)
        \end{cases}
    \end{gather}
    
    \columnbreak
    \begin{gather}
        f: A \cup B \to \mathbb{N}\\
        (A \cup B)_n \to n
    \end{gather}
\end{multicols}

This bijection establishes that $A \cup B$ is countable.

\subsection{Solution (b)}
This can be done using the winding bijection, detailed in the textbook.
The table for this is below.
\begin{center}
    \begin{tabular}{c| c c c c c}
            &1  &2  &3  &4  &\dots\\
        \hline
        $A_1$ &$a_{11}$    &$a_{21}$    &$a_{31}$   &$a_{41}$   &\dots\\
        $A_2$ &$a_{12}$    &$a_{22}$    &$a_{32}$   &$a_{42}$   &\dots\\
        $A_3$ &$a_{13}$    &$a_{23}$    &$a_{33}$   &$a_{43}$   &\dots\\
        $A_4$ &$a_{14}$    &$a_{24}$    &$a_{34}$   &$a_{44}$   &\dots\\
        \vdots  &\vdots    &\vdots      &\vdots     &\vdots     &$\ddots$
    \end{tabular}
\end{center}
This visual proof establishes the solution.

\pagebreak
\section{Exercise 2.6}
Show that $|\mathbb{N}| = |\mathbb{Z}|$ by finding an explicit bijection  from $\mathbb{N}$ to $\mathbb{Z}$. You do not need to prove your bijection works. 

\subsection*{Solution}
\begin{gather}
    f: \mathbb{N \to Z}\\
    n \to \lfloor\frac{n}{2}\rfloor (-1)^n
\end{gather}

\pagebreak
\section{Exercise 2.7}
Let $A, B \subseteq \mathbb{R}$. We define  $A \cdot B = \{a \cdot b: a \in A \andtxt b \in B\}$.  

(a) Give an example of sets $A_1$ and $B_1$ where $\left|A_1 \cdot B_1\right| < \max\{|A_1|, |B_1|\}$.  

(b) Give an example of sets $A_2$ and $B_2$ where $\left|A_2 \cdot B_2\right| > \max\{|A_2|, |B_2|\}$.

(c) Give an example of sets $A_3$ and $B_3$ where $\left|A_3 \cdot B_3\right| = \max\{|A_3|, |B_3|\}$.

(d) For which of the above does there exist an example where one or both of the sets are infinite? 

\subsection{Solution (a)}
\begin{gather}
    A_1 = \mathbb{N}\\
    B_1 = \{0\}\\
    A_1 \cdot B_1 = \{0\}
\end{gather}

\subsection{Solution (b)}
\begin{gather}
    A_2 = \{n : \mathbb{N} \andtxt 0 < n < 100\}\\
    B_2 = \{-1, 1\}\\
    A_2 \cdot B_2 = \{-99, -98, \dots , -1, 1, 2, \dots , 98, 99\}
\end{gather}

\subsection{Solution (c)}
\begin{gather}
    A_3 = \{1, 2, 3, 4, 5\}\\
    B_3 = \{1\}\\
    A_3 \cdot B_3 = \{1,2,3,4,5\}
\end{gather}

\subsection{Solution (d)}
For (a) and (c), yes.

\pagebreak
\section{Exercise 2.8}
(a) Describe a way to partition the set $\mathbb{N}$ into 6 subsets, each  containing infinitely many elements.\\ 
(b) Describe a way to partition the set $\mathbb{N}$ into infinitely many subsets, each containing infinitely many elements.

\subsection{Solution (a)}
Divide each number by 6 and take the remainder. 
The value of each remainder determines the set it would go in. The way to define each set is as follows:
\begin{align}
    &f: \mathbb{N} \to N_1  &n \to (n-1)*6 + 1\\
    &f: \mathbb{N} \to N_2  &n \to (n-1)*6 + 2\\
    &f: \mathbb{N} \to N_3  &n \to (n-1)*6 + 3\\
    &f: \mathbb{N} \to N_4  &n \to (n-1)*6 + 4\\
    &f: \mathbb{N} \to N_5  &n \to (n-1)*6 + 5\\
    &f: \mathbb{N} \to N_6  &n \to (n-1)*6 + 6
\end{align}

\subsection{Solution (b)}
Create one set $S_1$ and add one object: 1.
Create a second set $S_2$. 
Add 2 to $S_1$ and 3 to $S_2$.
Create $S_3$.
Add 4 to $S_1$, 5 to $S_2$, and 6 to $S_3$.
Your sets should look like the below.
\begin{align}
    S_1 &=  \{1,2,4\}\\
    S_2 &=  \{3,5\}\\
    S_3 &=  \{6\}
\end{align}

A pattern should be discernable here. 
We can define each individual set.
\begin{gather}
    f: \mathbb{N} \to S_i\\
    n \to \overunderset{i}{x = 0}{\Sigma} (x) + \overunderset{i + n - 1}{x = i}{\Sigma} (x) 
\end{gather}

\pagebreak
\section{Exercise 2.9}
Is $|\mathbb{Z \times N}|$ countable or uncountable? 

\subsection*{Solution}
Previously mentioned is the digonalization and pairing up practice, which has established $\mathbb{Q}$ as being countably infinite. 
In this instance, we can do the same, but instead of fractions, we can put $\mathbb{N}$ on one side and $\mathbb{Z}$ on the other. 
In the center will be the pairing up of the numbers.
This gives us the below table.
\begin{center}
    \begin{tabular}{c| c c c c c c}
            &0  &1  &-1 &2  &-2 &\dots\\
        \hline
        $1$ &\{0,1\}    &\{1,1\}    &\{-1,1\}   &\{2,1\}    &\{-2,1\}   &\dots\\
        $2$ &\{0,2\}    &\{1,2\}    &\{-1,2\}   &\{2,2\}    &\{-2,2\}   &\dots\\
        $3$ &\{0,3\}    &\{1,3\}    &\{-1,3\}   &\{2,3\}    &\{-2,3\}   &\dots\\
        $4$ &\{0,4\}    &\{1,4\}    &\{-1,4\}   &\{2,4\}    &\{-2,4\}   &\dots\\
        \vdots&\vdots   &\vdots     &\vdots     &\vdots     &\vdots     &$\ddots$
    \end{tabular}
\end{center}

We can subsequently use the winding bijection on this to get our countably infinite set.
As a note, this can be done in turn with any pair of countably infinite sets paired together in some way, as long as they are mapped in terms of some function $f(a,b)$.

\pagebreak
\section{Exercise 2.10}
Let $S$ be the set of sequences ($a_n$) where, for  each $n, a_n \in \{0, 1\}$. Is $S$ countable or uncountable? 

\subsection*{Solution}
Note: Each sequence is infinite in this case.

Cantor's proof that $\mathbb{R}$ is uncountably infinite may be applicable here.
Suppose that we laid out an array of all the sequences as binary decimals less than 1.
\begin{align}
    &0.a_{11}a_{12}a_{13}a_{14}a_{15}\\
    &0.a_{21}a_{22}a_{23}a_{24}a_{25}\\
    &0.a_{31}a_{32}a_{33}a_{34}a_{35}\\
    &0.a_{41}a_{42}a_{43}a_{44}a_{45}\\
    &0.a_{51}a_{52}a_{53}a_{54}a_{55}
\end{align}

Suppose that we were to create a member of this set.
We could take every value of each member of each set such that the first designating number matches the second designating number (e.g. $a_{11}$, $a_{22}$, $a_{33}$, $a_{xx}$, \dots). 
For each value of that, we would create a formula for the respective value of member $b_i$.
\begin{equation}
    a_{ii} \to b_i = \begin{cases}
        1   &\text{if } a_{ii} = 0\\
        0   &\text{if } a_{ii} = 1
    \end{cases}
\end{equation}

This in turn creates a new value of a set in $S$, which is previously unwritten.
This in turn establishes that the set is uncountably infinite.

\pagebreak
\section{Exercise 2.11}
Suppose that $X$ is a nonempty set. Prove that  the following three assertions are equivalent:  

(a) $X$ is finite or countably infinite.  

(b) There is a one-to-one function $f: X \to \mathbb{N}$.  

(c) There is an onto function $g: \mathbb{N} \to X$. 

\subsection*{Solution}
\begin{proof}
    We can rewrite (b) as for $a,b \in X$ and given $f(a), f(b) \in \mathbb{N}$, the below applies.
    \begin{equation}
        f(a) = f(b) \implies a = b 
    \end{equation}

    We can rewrite (c) as for $b \in X$, the below applies. 
    \begin{equation}
        \exists a \in \mathbb{N} : g(a) = b
    \end{equation}

    We can rewrite (a) as there being a correspondence between $\mathbb{N}$ and $X$ such that every .
\end{proof}

\pagebreak
\section{Exercise 2.12}
(a) Give an example of a collection of countably many disjoint open intervals, or prove that this does not exist. \\
(b) Give an example of a collection of uncountably many disjoint open intervals, or prove that this does not exist. 

\pagebreak
\section{Exercise 2.21}
Prove that the set of points on the unit circle in $\mathbb{R}^2$ (that is, $\{(x, y): x^2 + y^2 = 1\}$) is uncountable. 

\subsection*{Solution}
\begin{proof}
    As given in the definition of a unit circle, there is a formula that $(x,y)$ must follow.
    \begin{equation}
        x^2 + y^2 = 1
    \end{equation}
    
    This equation can be solved for $y$. 
    \begin{equation}
        y = \pm\sqrt{1 - x^2}
    \end{equation}
    
    This in turn creates a bijection between $\mathbb{R}$ and either $\mathbb{R} \times [0,\infty)$ or $\mathbb{R} \times (-\infty,0]$, for which we can focus on the former.
    This bijection can have an associated function.
    \begin{gather}
        f: \mathbb{R} \to \mathbb{R} \times [0,\infty)\\
        x \mapsto (x, \sqrt{1 - x^2})
    \end{gather}

    For this, there should be a correction that $\sqrt{1 - x^2}$ will only be real if $-1 \leq x \leq 1$, which we can use to correct our function.
    \begin{gather}
        f: [-1,1] \to [-1,1] \times [0,1]\\
        x \mapsto (x, \sqrt{1 - x^2})
    \end{gather}
    
    This gives us a domain for the points on a quarter of a unit circle of $[-1,1]$. 
    It has been previously established that there are uncountable real numbers between 0 and 1. 
    This means that there are uncountable real numbers between -1 and 1.
    As such, there are unlimited real numbers that are points on a unit circle.

    Q.E.D.
\end{proof}

\pagebreak
\section{Exercise 2.22}
A real number $x$ is said to be \textit{algebraic} (over the rationals) if it satisfies some polynomial equation (of positive degree) 
\begin{equation*}
    a_n x^n + a_{n-1} x^{n-1} + a_{n-2} x^{n-2} + \dots + a_1 x + a_0 = 0 
\end{equation*}
where each $a_i \in \mathbb{Q}$. If a real number is not algebraic, then it is \textit{transcendental}. 

(a) Prove that there are countably many algebraic numbers. (You  may use the fundamental theorem of algebra which says that a polynomial with degree $n$ has at most $n$ real roots.) 

(b) Prove that there are uncountably many transcendental numbers. 

\subsection{Solution (a)}
\begin{proof}
    The fundamental theorem of algebra says that a polynomial of degree $n$ has at most $n$ real roots.
    This means that the real roots of any given polynomial can be expressed as a finite set of real numbers and as such is countable.
    This means that any given polynomial has countably many real roots, all of which would be algebraic real numbers.

    It has been previously established that that the set of all polynomials with rational coefficients (as we are dealing with presently) is countable.
    We can in turn count the algebraic numbers by counting each polynomial and counting each of their roots as we count them.
    This establishes these algebraic numbers as countable.

    Q.E.D.
\end{proof}

\subsection{Solution (b)}
We can prove this in a similar way to proving there are uncountably many irrational numbers.
\begin{proof}[Proof by Contradiction]
    Suppose for contradication that the set of transcendental numbers $T$ is countable.
    The set of algebraic numbers $A$ is already established to be countable.
    We can express this in terms of cardinalities.
    \begin{equation}
        |A| = |T| = |\mathbb{N}|
    \end{equation}

    Transcendental numbers are defined as being real numbers that are not algebraic.
    \begin{gather}
        T = \mathbb{R} \setminus A\\
        A \cup T = \mathbb{R}
    \end{gather}

    Since both $A$ and $T$ are countable, we can form a function to combine the two through $\mathbb{N}$ which becomes $\mathbb{R}$ given $A \cup T = \mathbb{R}$.
    \begin{gather}
        f: \mathbb{N} \to \mathbb{R}\\
        n \mapsto \begin{cases}
            A(\lceil\frac{n}{2}\rceil) &\text{if } n \equiv 0 \mod 2\\
            T(\lceil\frac{n}{2}\rceil) &\text{if } n \equiv 1 \mod 2
        \end{cases}
    \end{gather}

    This establishes that $\mathbb{R}$ is countably infinite. 
    This contradicts that $\mathbb{R}$ is uncountably infinite.
    As such, we can conclude that there are uncountably many transcendental numbers.
\end{proof}

\end{document}